\clearpage
\section{Assembly Language and Canonical Text}

This section defines the canonical textual mapping for one decoded instruction record. The syntax lines in the
instruction entries are form patterns: field markers in parentheses identify encoded fields, \texttt{X(z:...)}
lists the size suffixes selected by field \texttt{z}, and angle-bracketed names identify appended payloads. Those
markers explain the encoding and are replaced by concrete operands in assembly source.

\subsection{Canonical Statement Grammar}

The grammar consumes one instruction statement after symbol and relocation expressions have been resolved to the
integer value presented to instruction-form selection.

\begin{manualcode}
\manualcodeline{statement := [ LEN integer , ] instruction}
\manualcodeline{instruction := instruction-head [ operand-list ]}
\manualcodeline{instruction-head := base-mnemonic [ condition-name ] [ size-suffix ] [ / memory-order ]}
\manualcodeline{operand-list := operand \{ , operand \}}
\manualcodeline{operand := register | immediate | effective-address | flag-mask | instruction-body}
\manualcodeline{register := R0 ... R15 | SP | PC | F0 ... F15 | architectural named register}
\manualcodeline{integer := [ + | - ] ( decimal | 0x hexadecimal | 0b binary )}
\end{manualcode}

Canonical mnemonics, register names, segment-register names, condition names, size suffixes, memory-order names,
and flag names use uppercase spelling. Decimal digits have no prefix. Hexadecimal and binary literals use the
lowercase prefixes \texttt{0x} and \texttt{0b}; canonical hexadecimal digits are lowercase. Canonical signed
immediates and displacements use decimal notation, while absolute addresses use hexadecimal notation.

For a conditional mnemonic, the primary condition name from the Condition Encoding table immediately follows the
base mnemonic. A size suffix follows the condition. Thus a conditional long form is written
\texttt{JNE.L}, not with the \texttt{cc}, \texttt{X}, \texttt{z}, or field-marker notation used by an encoding
pattern. An atomic memory-order selector follows the size suffix, as in
\texttt{CMPXCHG.Q/ACQUIRE}.

Operands appear in the order printed by the selected instruction form. A two-operand data form prints source,
comma, destination. A three-operand form retains the names and order stated by its instruction entry. Canonical
output places one space after a mnemonic head and after each comma, and one space on each side of an address
\texttt{+} or \texttt{-} operator.

\subsection{Effective-Address Text}

The canonical effective-address productions are:

\begin{manualcode}
\manualcodeline{effective-address := register | immediate | memory-address}
\manualcodeline{memory-address := [ address-expression ]}
\manualcodeline{address-expression := base | base + displacement | base - displacement}
\manualcodeline{address-expression := segment : base [ + index * scale ] [ + displacement | - displacement ]}
\manualcodeline{address-expression := base + index * scale [ + displacement | - displacement ]}
\manualcodeline{base := R0 ... R15 | SP | PC | 0 | R0++ ... R15++ | {-}{-}R0 ... {-}{-}R15}
\manualcodeline{index := R0 ... R15 | R0++ ... R15++ | {-}{-}R0 ... {-}{-}R15}
\manualcodeline{segment := DS | SS | GS0 ... GS5}
\manualcodeline{scale := 1 | 2 | 4 | 8}
\end{manualcode}

The selected compact or EXT0 profile determines which production members are valid for that instruction field.
The register, descriptor, displacement, and scale values in the text map directly to the fields and payload order
specified in the Effective Addressing chapters. A canonical disassembler chooses the primary condition name and
the concrete register and segment names from their encoding tables.

\subsection{Aliases and Canonical Disassembly}

\manualtablecaption{Assembler Aliases and Canonical Text}
\begin{manualtabular}{@{}>{\raggedright\arraybackslash}p{1.35in}>{\raggedright\arraybackslash}p{1.55in}>{\raggedright\arraybackslash}p{2.55in}@{}}
\toprule
\rowcolor{ManualHeaderFill}
\textbf{Source spelling} & \textbf{Canonical text} & \textbf{Mapping}\\
\midrule
\texttt{CLC} & \texttt{CLRF C} & Encodes the C-bit flag mask.\\
\texttt{STC} & \texttt{SETF C} & Encodes the C-bit flag mask.\\
\texttt{REP Rn, (instruction)} & \texttt{REP Rn, (instruction)} & Encodes condition T in the REPcc form.\\
\texttt{REPGF Rn, \{ body \}} & \texttt{REPG Rn, body\_bytes} & Applies the REPGF assembly-time body check, then emits REPG.\\
condition alias & primary condition name & Uses the primary name in the Condition Encoding table.\\
\bottomrule
\end{manualtabular}

\texttt{LEN n, instruction} selects an extended record of exactly \texttt{n} bytes. Form selection first chooses
the shortest concrete form whose fields, operand classes, ranges, and EA constraints match the statement. The
explicit length then supplies zero padding through byte \texttt{n-1}. Without \texttt{LEN}, the selected concrete
form uses its shortest valid record. Canonical disassembly emits \texttt{LEN n,} exactly when an extended record is
longer than its concrete form, as defined in Instruction Header Formats.

\subsection{Form Selection and Diagnostics}

Form matching evaluates the complete operand list, explicit suffixes and qualifiers, field ranges, EA profile
constraints, destination-overlap rules, and the requested record length. The assembler reports the stage that
prevented a unique encoding:

\begin{itemize}
\item \textbf{lexical}: a token does not match the statement grammar;
\item \textbf{expression}: an operand expression has no resolved integer or relocation value;
\item \textbf{form}: no instruction form has the supplied mnemonic head and operand classes;
\item \textbf{range}: a resolved value does not fit the selected field or payload width;
\item \textbf{constraint}: an EA profile, selector value, or operand relation fails the selected form's rule;
\item \textbf{length}: a \texttt{LEN} value lies outside 3 through 18 or cannot contain the concrete form; and
\item \textbf{ambiguity}: more than one shortest concrete form remains after all preceding checks.
\end{itemize}

\subsection{Conformance Encodings}

The complete source cases are recorded in
\texttt{isa/\allowbreak{}reference/\allowbreak{}assembler\_\allowbreak{}conformance\_\allowbreak{}vectors.yaml}.
The following cases show the byte order used by that
source. Fields are listed after decoding the header and opcode payload; appended descriptor, immediate,
displacement, and padding bytes follow in instruction-stream order.

\manualtablecaption{Canonical Assembly Conformance Encodings}
\begin{manuallongtable}{@{}>{\raggedright\arraybackslash}p{1.35in}>{\raggedright\arraybackslash}p{1.65in}>{\raggedright\arraybackslash}p{1.10in}>{\raggedright\arraybackslash}p{1.45in}@{}}
\toprule
\rowcolor{ManualHeaderFill}
\textbf{Assembly} & \textbf{Form and fields} & \textbf{Bytes} & \textbf{Canonical disassembly}\\
\midrule
\endfirsthead
\multicolumn{4}{l}{\scriptsize\itshape Table \themanualtable\ (continued)}\\
\toprule
\rowcolor{ManualHeaderFill}
\textbf{Assembly} & \textbf{Form and fields} & \textbf{Bytes} & \textbf{Canonical disassembly}\\
\midrule
\endhead
@ASSEMBLER_CONFORMANCE_ROWS@
\bottomrule
\end{manuallongtable}
